\documentclass[12pt,a4paper]{article}

\usepackage[T1]{fontenc}
\usepackage[utf8]{inputenc}
\usepackage[portuguese]{babel}
\usepackage{geometry}
\geometry{margin=2.5cm}

\usepackage{amsmath,amssymb}
\usepackage{booktabs}
\usepackage{array}
\usepackage{longtable}
\usepackage{multirow}
\usepackage{graphicx}
\usepackage{xcolor}
\usepackage{hyperref}
\usepackage{float}
\usepackage{caption}
\usepackage{subcaption}
\usepackage{enumitem}
\usepackage{tikz}
\usepackage[edges]{forest}
\usetikzlibrary{arrows.meta,positioning,calc}

\usepackage{listings}
\lstdefinestyle{py}{
  language=Python,
  basicstyle=\ttfamily\small,
  keywordstyle=\color{blue!70!black},
  commentstyle=\color{green!40!black},
  stringstyle=\color{orange!60!black},
  numbers=left,
  numberstyle=\scriptsize\color{gray},
  stepnumber=1,
  numbersep=8pt,
  showstringspaces=false,
  breaklines=true,
  frame=single,
  tabsize=4,
  extendedchars=true,
  literate={á}{{\'a}}1 {ã}{{\~a}}1 {â}{{\^a}}1 {à}{{\`a}}1 {é}{{\'e}}1 {ê}{{\^e}}1 {í}{{\'i}}1 {ó}{{\'o}}1 {õ}{{\~o}}1 {ô}{{\^o}}1 {ú}{{\'u}}1 {ç}{{\c{c}}}1 {Á}{{\'A}}1 {Ã}{{\~A}}1 {Â}{{\^A}}1 {À}{{\`A}}1 {É}{{\'E}}1 {Ê}{{\^E}}1 {Í}{{\'I}}1 {Ó}{{\'O}}1 {Õ}{{\~O}}1 {Ô}{{\^O}}1 {Ú}{{\'U}}1 {Ç}{{\c{C}}}1
}
\lstset{style=py}

\newcommand{\estado}[1]{\texttt{#1}}
\newcommand{\fronteira}{\textit{fronteira}}

\title{Prova 1 - Inteligência Artificial\\Resolução de Questões}
\author{Aluno: Diego}
\date{Maio de 2026}

\begin{document}

\maketitle
\tableofcontents
\newpage

\section*{Orientações}
\begin{itemize}[leftmargin=1.5em]
  \item Resolução detalhada das questões da Prova 1 de IA.
  \item As árvores de busca foram geradas utilizando \texttt{tikz}.
  \item Os códigos Python utilizados para os experimentos estão anexados no apêndice.
\end{itemize}

\section*{Resolução da Questão 1: Algoritmos de Busca em Grafos}

\subsection*{A) Estruturação do Problema}

A definição deste desafio de busca em espaço de estados baseia-se nos seguintes elementos essenciais:

\begin{itemize}
    \item \textbf{Representação do Estado Atual}: Uma letra ou string que identifica o aposento atual do agente. Por exemplo: $s \in \{A, B, C, \dots, S\}$.
    \item \textbf{Estado de Origem}: O ponto de partida do labirinto é a sala $A$, portanto $s_0 = A$.
    \item \textbf{Verificação de Conclusão}: Condição que testa se a sala atual $s$ corresponde ao alvo $S$. A função $ChegouAoFim(s)$ retorna $Verdadeiro$ caso $s = S$, e $Falso$ em outras situações.
    \item \textbf{Transições Possíveis (Sucessores)}: Função que mapeia os deslocamentos válidos a partir de uma sala $s$ para as adjacentes, seguindo rigidamente a ordem alfabética. Ex: $Vizinhos(A) = \{B, C, D\}$, $Vizinhos(B) = \{E, F\}$.
    \item \textbf{Métrica de Custo}: O custo associado à transição entre duas salas conectadas. Sendo um movimento de custo uniforme, $c(s, a, s') = 1$. O custo total é a quantidade de passos dados.
\end{itemize}

\textbf{Definição do Grafo em Python:}
\begin{lstlisting}[language=Python]
adjacency_map = {
    'A': ['B', 'C', 'D'], 'B': ['E', 'F'], 'C': ['G', 'H'],
    'D': ['I'], 'E': ['J'], 'F': ['K', 'L'], 'G': ['M'],
    'H': ['N', 'O'], 'I': ['P'], 'J': [], 'K': ['Q'],
    'L': [], 'M': ['R'], 'N': [], 'O': ['S'], 'P': [],
    'Q': [], 'R': [], 'S': []
}
start_node = 'A'
target_node = 'S'
\end{lstlisting}

\subsection*{B) Avaliação dos Métodos de Busca}

\subsubsection*{1. Busca em Largura (BFS)}

\textbf{Desempenho da BFS:}
\begin{itemize}
    \item \textbf{Trajeto da Solução:} A $\rightarrow$ C $\rightarrow$ H $\rightarrow$ O $\rightarrow$ S
    \item \textbf{Custo Total:} 4
    \item \textbf{Nível de Profundidade:} 4
    \item \textbf{Quantidade de Nós Expandidos:} 15
    \item \textbf{Quantidade de Nós Gerados:} 19
\end{itemize}

\textbf{Acompanhamento da Execução (Fronteira FIFO):}
\begin{table}[h!]
\centering
\begin{tabular}{|c|c|l|}
\hline
\textbf{Iteração} & \textbf{Nó Processado} & \textbf{Estado da Fronteira Pós-Expansão} \\ \hline
1 & A & [B, C, D] \\ \hline
2 & B & [C, D, E, F] \\ \hline
3 & C & [D, E, F, G, H] \\ \hline
4 & D & [E, F, G, H, I] \\ \hline
5 & E & [F, G, H, I, J] \\ \hline
6 & F & [G, H, I, J, K, L] \\ \hline
7 & G & [H, I, J, K, L, M] \\ \hline
8 & H & [I, J, K, L, M, N, O] \\ \hline
9 & I & [J, K, L, M, N, O, P] \\ \hline
10 & J & [K, L, M, N, O, P] \\ \hline
11 & K & [L, M, N, O, P, Q] \\ \hline
12 & L & [M, N, O, P, Q] \\ \hline
13 & M & [N, O, P, Q, R] \\ \hline
14 & N & [O, P, Q, R] \\ \hline
15 & O & [P, Q, R, S] \\ \hline
16 & S & Sucesso (Objetivo Encontrado) \\ \hline
\end{tabular}
\end{table}

\textbf{Árvore de Exploração Parcial (BFS):}
\textit{Nota: Nós preenchidos em ciano (\textcolor{cyan}{Cor}) indicam expansão. O alvo é laranja (\textcolor{orange}{Cor}). Círculos vazios apenas gerados.}

\begin{center}
\begin{tikzpicture}[level distance=1.2cm,
  level 1/.style={sibling distance=4cm},
  level 2/.style={sibling distance=2cm},
  level 3/.style={sibling distance=1.5cm},
  every node/.style={circle, draw, minimum size=0.6cm, inner sep=1pt}]
  
  \node[fill=cyan!20] {A}
    child {node[fill=cyan!20] {B}
      child {node[fill=cyan!20] {E}
        child {node[fill=cyan!20] {J}}
      }
      child {node[fill=cyan!20] {F}
        child {node[fill=cyan!20] {K}
          child {node {Q}}
        }
        child {node[fill=cyan!20] {L}}
      }
    }
    child {node[fill=cyan!20] {C}
      child {node[fill=cyan!20] {G}
        child {node[fill=cyan!20] {M}
          child {node {R}}
        }
      }
      child {node[fill=cyan!20] {H}
        child {node[fill=cyan!20] {N}}
        child {node[fill=cyan!20] {O}
          child {node[fill=orange!40] {S}}
        }
      }
    }
    child {node[fill=cyan!20] {D}
      child {node[fill=cyan!20] {I}
        child {node {P}}
      }
    };
\end{tikzpicture}
\end{center}

\subsubsection*{2. Busca em Profundidade (DFS)}

Para garantir que a DFS priorize os vizinhos iniciais (por exemplo, visitar B antes de C ou D), a fronteira do tipo Pilha insere os elementos em ordem reversa: D, depois C, depois B.

\textbf{Desempenho da DFS:}
\begin{itemize}
    \item \textbf{Trajeto da Solução:} A $\rightarrow$ C $\rightarrow$ H $\rightarrow$ O $\rightarrow$ S
    \item \textbf{Custo Total:} 4
    \item \textbf{Nível de Profundidade:} 4
    \item \textbf{Quantidade de Nós Expandidos:} 15
    \item \textbf{Quantidade de Nós Gerados:} 17
\end{itemize}

\textbf{Acompanhamento da Execução (Fronteira LIFO):}
\begin{table}[h!]
\centering
\begin{tabular}{|c|c|l|}
\hline
\textbf{Iteração} & \textbf{Nó Processado} & \textbf{Estado da Fronteira Pós-Expansão} \\ \hline
1 & A & [D, C, B] \\ \hline
2 & B & [D, C, F, E] \\ \hline
3 & E & [D, C, F, J] \\ \hline
4 & J & [D, C, F] \\ \hline
5 & F & [D, C, L, K] \\ \hline
6 & K & [D, C, L, Q] \\ \hline
7 & Q & [D, C, L] \\ \hline
8 & L & [D, C] \\ \hline
9 & C & [D, H, G] \\ \hline
10 & G & [D, H, M] \\ \hline
11 & M & [D, H, R] \\ \hline
12 & R & [D, H] \\ \hline
13 & H & [D, O, N] \\ \hline
14 & N & [D, O] \\ \hline
15 & O & [D, S] \\ \hline
16 & S & Sucesso (Objetivo Encontrado) \\ \hline
\end{tabular}
\end{table}

\textbf{Árvore de Exploração Parcial (DFS):}

\begin{center}
\begin{tikzpicture}[level distance=1.2cm,
  level 1/.style={sibling distance=4cm},
  level 2/.style={sibling distance=2cm},
  level 3/.style={sibling distance=1.5cm},
  every node/.style={circle, draw, minimum size=0.6cm, inner sep=1pt}]
  
  \node[fill=cyan!20] {A}
    child {node[fill=cyan!20] {B}
      child {node[fill=cyan!20] {E}
        child {node[fill=cyan!20] {J}}
      }
      child {node[fill=cyan!20] {F}
        child {node[fill=cyan!20] {K}
          child {node[fill=cyan!20] {Q}}
        }
        child {node[fill=cyan!20] {L}}
      }
    }
    child {node[fill=cyan!20] {C}
      child {node[fill=cyan!20] {G}
        child {node[fill=cyan!20] {M}
          child {node[fill=cyan!20] {R}}
        }
      }
      child {node[fill=cyan!20] {H}
        child {node[fill=cyan!20] {N}}
        child {node[fill=cyan!20] {O}
          child {node[fill=orange!40] {S}}
        }
      }
    }
    child {node {D}};
\end{tikzpicture}
\end{center}

\subsubsection*{3. Busca de Aprofundamento Iterativo (IDS)}

A estratégia IDS repete sondagens em profundidade limitada (DLS) nos limiares de $L=0$ até $L=4$, onde o objetivo $S$ é enfim descoberto.

\textbf{Desempenho da IDS:}
\begin{itemize}
    \item \textbf{Trajeto da Solução:} A $\rightarrow$ C $\rightarrow$ H $\rightarrow$ O $\rightarrow$ S
    \item \textbf{Custo Total:} 4
    \item \textbf{Nós Expandidos (Soma Global):} 27
    \item \textbf{Nós Gerados (Soma Global):} 47
\end{itemize}

O progresso repete a lógica até o limite 4. A árvore gerada no último limite ($L=4$) é idêntica à porção da árvore de DFS que alcança profundidade 4, resultando na mesma tabela observada para DFS.

\subsection*{C) Algoritmos em Linguagem Python}

O código-fonte com as implementações encontra-se na seção de Anexos do trabalho.

\subsection*{D) Experimento: Ordem de Geração Invertida C $\rightarrow$ H, G}

Ao modificar estritamente a prioridade dos vizinhos de C de \texttt{(G, H)} para \texttt{(H, G)}, o comportamento de percurso se ajusta. 

\subsubsection*{1. Nova Execução de BFS}
\begin{table}[h!]
\centering
\begin{tabular}{|c|c|l|}
\hline
\textbf{Iteração} & \textbf{Nó Processado} & \textbf{Estado da Fronteira Pós-Expansão} \\ \hline
1 & A & [B, C, D] \\ \hline
2 & B & [C, D, E, F] \\ \hline
3 & C & [D, E, F, H, G] \\ \hline
4 & D & [E, F, H, G, I] \\ \hline
5 & E & [F, H, G, I, J] \\ \hline
6 & F & [H, G, I, J, K, L] \\ \hline
7 & H & [G, I, J, K, L, N, O] \\ \hline
8 & G & [I, J, K, L, N, O, M] \\ \hline
9 & I & [J, K, L, N, O, M, P] \\ \hline
10 & J & [K, L, N, O, M, P] \\ \hline
11 & K & [L, N, O, M, P] \\ \hline
12 & L & [N, O, M, P, Q] \\ \hline
13 & N & [O, M, P, Q] \\ \hline
14 & O & [M, P, Q, S] \\ \hline
15 & S & Sucesso (Objetivo Encontrado) \\ \hline
\end{tabular}
\end{table}

\begin{center}
\begin{tikzpicture}[level distance=1.2cm, level 1/.style={sibling distance=4cm}, level 2/.style={sibling distance=2cm}, level 3/.style={sibling distance=1.5cm},
  every node/.style={circle, draw, minimum size=0.6cm, inner sep=1pt}]
  \node[fill=cyan!20] {A}
    child {node[fill=cyan!20] {B}
      child {node[fill=cyan!20] {E} child {node[fill=cyan!20] {J}}}
      child {node[fill=cyan!20] {F} child {node[fill=cyan!20] {K}} child {node[fill=cyan!20] {L}}}
    }
    child {node[fill=cyan!20] {C}
      child {node[fill=cyan!20] {H} child {node[fill=cyan!20] {N}} child {node[fill=cyan!20] {O} child {node[fill=orange!40] {S}}}}
      child {node[fill=cyan!20] {G} child {node {M}}}
    }
    child {node[fill=cyan!20] {D}
      child {node[fill=cyan!20] {I} child {node {P}}}
    };
\end{tikzpicture}
\end{center}

\subsubsection*{2. Nova Execução de DFS}
\begin{table}[h!]
\centering
\begin{tabular}{|c|c|l|}
\hline
\textbf{Iteração} & \textbf{Nó Processado} & \textbf{Estado da Fronteira Pós-Expansão} \\ \hline
1 & A & [D, C, B] \\ \hline
2 & B & [D, C, F, E] \\ \hline
3 & E & [D, C, F, J] \\ \hline
4 & J & [D, C, F] \\ \hline
5 & F & [D, C, L, K] \\ \hline
6 & K & [D, C, L, Q] \\ \hline
7 & Q & [D, C, L] \\ \hline
8 & L & [D, C] \\ \hline
9 & C & [D, G, H] \\ \hline
10 & H & [D, G, O, N] \\ \hline
11 & N & [D, G, O] \\ \hline
12 & O & [D, G, S] \\ \hline
13 & S & Sucesso (Objetivo Encontrado) \\ \hline
\end{tabular}
\end{table}

\begin{center}
\begin{tikzpicture}[level distance=1.2cm, level 1/.style={sibling distance=4cm}, level 2/.style={sibling distance=2cm}, level 3/.style={sibling distance=1.5cm},
  every node/.style={circle, draw, minimum size=0.6cm, inner sep=1pt}]
  \node[fill=cyan!20] {A}
    child {node[fill=cyan!20] {B}
      child {node[fill=cyan!20] {E} child {node[fill=cyan!20] {J}}}
      child {node[fill=cyan!20] {F} child {node[fill=cyan!20] {K} child {node[fill=cyan!20] {Q}}} child {node[fill=cyan!20] {L}}}
    }
    child {node[fill=cyan!20] {C}
      child {node[fill=cyan!20] {H} child {node[fill=cyan!20] {N}} child {node[fill=cyan!20] {O} child {node[fill=orange!40] {S}}}}
      child {node {G}}
    }
    child {node {D}};
\end{tikzpicture}
\end{center}

\subsubsection*{Conclusões do Experimento}
\begin{enumerate}[label=\roman*)]
    \item O percurso ótimo ($A \rightarrow C \rightarrow H \rightarrow O \rightarrow S$) foi mantido idêntico em todas as variantes.
    \item Ocorreu uma otimização nos nós expandidos em todas as técnicas.
    \item A Busca em Profundidade sofreu impacto tremendo (reduzindo a exploração inútil do ramo G de forma precoce).
    \item A Amplitude manteve-se como a técnica mais punitiva em RAM devido à retenção dos irmãos não-visitados de nível.
\end{enumerate}

\subsection*{E) Avaliação e Comparativo}

A DFS demonstrou forte dependência de ordenação arbitrária para sua eficiência de curto prazo, enquanto a BFS é mecanicamente imune a ordenações em termos de garantia de profundidade mínima, mas extremamente custosa em cenários reais com amplo fator de ramificação. A IDS desponta como a escolha teórica mais robusta na preservação de RAM com exploração sistêmica controlada.

\section*{Resolução da Questão 2: Busca Gulosa e A*}

\subsection*{A) Estruturação Lógica do Problema}

A formulação heurística do robô em malha com bloqueios:

\begin{itemize}
    \item \textbf{Formato de Estado}: Cada localização mapeada no piso é definida por um símbolo. Exemplo: $s \in \{A, B, C, \dots, T\}$.
    \item \textbf{Início}: A partida ocorre na célula $A$.
    \item \textbf{Encerramento}: Chega-se ao fim se e somente se $s = T$.
    \item \textbf{Regras de Transição}: Vizinhos disponíveis e seus pedágios de travessia. Por exemplo: $Vizinhos(A) = \{(B, 2), (C, 4), (D, 3)\}$.
    \item \textbf{Cálculo de Custos ($g$)}: O valor $g(n)$ rastreia o peso acumulado de todas as rotas percorridas desde o nó base $A$ até o destino local $n$.
    \item \textbf{Impacto da Heurística ($h$)}: Representação de um "palpite instruído" do esforço restante de $n$ até $T$. No A*, a subestimação estrita do verdadeiro custo ($h(n) \le h^*(n)$) confere à heurística o status de "admissível", travando a rota como matematicamente perfeita (ótima).
\end{itemize}

\subsection*{B) Avaliação com Greedy Best-First Search}

A tática gulosa guia as decisões avaliando cegamente $f(n) = h(n)$, ignorando pedágios já pagos.

\textbf{Desempenho (Greedy):}
\begin{itemize}
    \item \textbf{Trilha Encontrada:} A $\rightarrow$ C $\rightarrow$ H $\rightarrow$ O $\rightarrow$ T
    \item \textbf{Custo Acumulado:} 19
    \item \textbf{Quantia de Nós Expandidos:} 5
    \item \textbf{Quantia de Nós Gerados:} 9
\end{itemize}

\textbf{Passo-a-Passo Guloso:}
\begin{table}[h!]
\centering
\small
\begin{tabular}{|c|c|c|l|}
\hline
\textbf{Turno} & \textbf{Expansão} & \textbf{$h(n)$} & \textbf{Fila Prioritária por $h$} \\ \hline
1 & A & 10 & [C(7), B(8), D(9)] \\ \hline
2 & C & 7 & [H(4), G(6), B(8), D(9)] \\ \hline
3 & H & 4 & [O(1), N(4), G(6), B(8), D(9)] \\ \hline
4 & O & 1 & [T(0), N(4), G(6), B(8), D(9)] \\ \hline
5 & T & 0 & Concluído com sucesso \\ \hline
\end{tabular}
\end{table}

\textbf{Árvore de Exploração Parcial (Greedy):}
\textit{Preenchimento roxo (\textcolor{purple}{Cor}) para explorados. Verde (\textcolor{teal}{Cor}) para o nó final.}

\begin{center}
\begin{tikzpicture}[level distance=1.5cm,
  level 1/.style={sibling distance=3.5cm},
  level 2/.style={sibling distance=1.5cm},
  level 3/.style={sibling distance=1.2cm},
  every node/.style={circle, draw, minimum size=0.8cm, inner sep=1pt}]
  
  \node[fill=purple!20, label=right:{\scriptsize 10}] {A}
    child {node[label=left:{\scriptsize 8}] {B}}
    child {node[fill=purple!20, label=left:{\scriptsize 7}] {C}
        child {node[label=left:{\scriptsize 6}] {G}}
        child {node[fill=purple!20, label=right:{\scriptsize 4}] {H}
            child {node[label=left:{\scriptsize 4}] {N}}
            child {node[fill=purple!20, label=right:{\scriptsize 1}] {O}
                child {node[fill=teal!40, label=right:{\scriptsize 0}] {T}}
            }
        }
    }
    child {node[label=right:{\scriptsize 9}] {D}};
\end{tikzpicture}
\end{center}

\newpage
\subsection*{C) Avaliação com Algoritmo A*}

A* equilibra os pesos do passado e previsões futuras: $f(n) = g(n) + h(n)$.

\textbf{Desempenho (A*):}
\begin{itemize}
    \item \textbf{Trilha Encontrada:} A $\rightarrow$ B $\rightarrow$ F $\rightarrow$ K $\rightarrow$ R $\rightarrow$ T
    \item \textbf{Custo Acumulado:} 17
    \item \textbf{Quantia de Nós Expandidos:} 20
\end{itemize}

\textbf{Passo-a-Passo A*:}
\begin{table}[h!]
\centering
\scriptsize
\begin{tabular}{|c|c|c|c|c|l|l|}
\hline
\textbf{Trn} & \textbf{Nó} & \textbf{$g$} & \textbf{$h$} & \textbf{$f$} & \textbf{Fila (Nó:f)} & \textbf{Rota} \\ \hline
1 & A & 0 & 10 & 10 & [B:10, C:11, D:12] & A \\ \hline
2 & B & 2 & 8 & 10 & [C:11, E:11, D:12, F:12] & A-B \\ \hline
3 & C & 4 & 7 & 11 & [E:11, D:12, F:12, G:14, H:14] & A-C \\ \hline
4 & E & 5 & 6 & 11 & [D:12, F:12, G:14, H:14, J:14] & A-B-E \\ \hline
5 & D & 3 & 9 & 12 & [F:12, I:12, G:14, H:14, J:14] & A-D \\ \hline
6 & F & 7 & 5 & 12 & [I:12, K:13, G:14, H:14, J:14, L:18] & A-B-F \\ \hline
7 & I & 5 & 7 & 12 & [K:13, G:14, H:14, J:14, L:18, P:18] & A-D-I \\ \hline
8 & K & 10 & 3 & 13 & [G:14, H:14, J:14, R:15, L:18, P:18] & A-B-F-K \\ \hline
9 & G & 8 & 6 & 14 & [H:14, J:14, R:15, M:17, L:18, P:18] & A-C-G \\ \hline
10 & H & 10 & 4 & 14 & [J:14, O:15, R:15, M:17, N:17, L:18, P:18] & A-C-H \\ \hline
11 & J & 9 & 5 & 14 & [O:15, R:15, M:17, N:17, Q:17, L:18, P:18] & A-B-E-J \\ \hline
12 & O & 14 & 1 & 15 & [R:15, M:17, N:17, Q:17, L:18, P:18, T:19] & A-C-H-O \\ \hline
13 & R & 13 & 2 & 15 & [M:17, N:17, Q:17, T:17, L:18, P:18] & A-B-F-K-R \\ \hline
14 & M & 14 & 3 & 17 & [N:17, Q:17, S:17, T:17, L:18, P:18] & A-C-G-M \\ \hline
15 & N & 13 & 4 & 17 & [Q:17, S:17, T:17, L:18, P:18] & A-C-H-N \\ \hline
16 & Q & 13 & 4 & 17 & [S:17, T:17, L:18, P:18] & A-B-E-J-Q \\ \hline
17 & S & 16 & 1 & 17 & [T:17, L:18, P:18] & A-C-G-M-S \\ \hline
18 & L & 12 & 6 & 18 & [T:17, P:18] & A-B-F-L \\ \hline
19 & P & 10 & 8 & 18 & [T:17] & A-D-I-P \\ \hline
20 & T & 17 & 0 & 17 & Sucesso & A-B-F-K-R-T \\ \hline
\end{tabular}
\end{table}

\textbf{Árvore de Exploração Parcial (A*):}
\textit{Legenda: Expansões em roxo (\textcolor{purple}{Cor}). Trilha perfeita em negrito com tom verde-escuro.}

\begin{center}
\begin{tikzpicture}[level distance=1.3cm,
  level 1/.style={sibling distance=4cm},
  level 2/.style={sibling distance=2cm},
  level 3/.style={sibling distance=1cm},
  every node/.style={circle, draw, minimum size=0.7cm, inner sep=1pt}]
  
  \node[fill=purple!20, font=\bfseries\color{teal!60!black}] {A}
    child {node[fill=purple!20, font=\bfseries\color{teal!60!black}] {B}
      child {node[fill=purple!20] {E} child {node[fill=purple!20] {J} child {node[fill=purple!20] {Q}}}}
      child {node[fill=purple!20, font=\bfseries\color{teal!60!black}] {F} 
        child {node[fill=purple!20, font=\bfseries\color{teal!60!black}] {K} child {node[fill=purple!20, font=\bfseries\color{teal!60!black}] {R} child {node[fill=teal!40] {T}}}}
        child {node[fill=purple!20] {L}}
      }
    }
    child {node[fill=purple!20] {C}
      child {node[fill=purple!20] {G} child {node[fill=purple!20] {M} child {node[fill=purple!20] {S} child {node {T}}}}}
      child {node[fill=purple!20] {H} child {node[fill=purple!20] {N}} child {node[fill=purple!20] {O} child {node {T}}}}
    }
    child {node[fill=purple!20] {D}
      child {node[fill=purple!20] {I} child {node[fill=purple!20] {P}}}
    };
\end{tikzpicture}
\end{center}

\subsection*{D) Experimento: Distorção das Heurísticas}

Injetamos erros nos sensores de estimativa:
\begin{itemize}
    \item $h(D)$ alterado para $2$
    \item $h(I)$ alterado para $1$
    \item $h(B)$ alterado de 8 para $22$
\end{itemize}

\textbf{Comparativo do Impacto:}
\begin{table}[h!]
\centering
\begin{tabular}{|l|l|c|c|}
\hline
\textbf{Abordagem} & \textbf{Rota Consolidada} & \textbf{Preço} & \textbf{Expansões} \\ \hline
Guloso (Padrão) & A $\rightarrow$ C $\rightarrow$ H $\rightarrow$ O $\rightarrow$ T & 19 & 5 \\ \hline
Guloso (Corrompido) & A $\rightarrow$ C $\rightarrow$ H $\rightarrow$ O $\rightarrow$ T & 19 & 7 \\ \hline
A* (Padrão) & A $\rightarrow$ B $\rightarrow$ F $\rightarrow$ K $\rightarrow$ R $\rightarrow$ T & 17 & 20 \\ \hline
A* (Corrompido) & A $\rightarrow$ C $\rightarrow$ H $\rightarrow$ O $\rightarrow$ T & 19 & 13 \\ \hline
\end{tabular}
\end{table}

\textbf{Síntese dos Efeitos:}
\begin{enumerate}[label=\roman*)]
    \item A rota A* degenerou, abandonando o caminho ideal (17) por um sub-ótimo (19).
    \item A causa direta da falha do A* reside no fato de $h(B)=22$ ser uma superestimativa extrema, violando a prova matemática de admissibilidade do A*.
    \item A pesquisa Gulosa reagiu vagamente nos fundos falsos (testou D primeiro pois $h=2$), mas em geral re-encontrou a mesma trilha sub-ótima sem maiores prejuízos lógicos.
\end{enumerate}

\subsection*{E) Conclusões Gerais}

O A* é incontestável como líder em confiabilidade para malhas difíceis. Enquanto a busca gulosa desce rapidamente rampa abaixo (frequentemente trombando em obstáculos e achando percursos longos em forma de $U$), o A* se resguarda conferindo $g(n)$, e é o único blindado com garantias matemáticas de otimalidade na literatura, condicionada à pureza de seu estimador heurístico $h(n)$.

\section*{Resolução da Questão 3: Busca Local (N-Rainhas)}

\subsection*{A) Estruturação Computacional do Enigma}

A missão é depositar 8 rainhas sobre um tabuleiro de xadrez $8 \times 8$ com zero colisões (ataques cruzados). Traduzido para um algoritmo de Hill-Climbing, temos:

\begin{itemize}
    \item \textbf{Codificação (Estado):} Array $S$ de 8 inteiros usando indexação base-zero (0 a 7). O valor em $S[c]$ indica a linha em que a rainha da coluna $c$ foi inserida. Ex: $[0, 0, 0, 0, 0, 0, 0, 0]$ (todas empilhadas na primeira linha).
    \item \textbf{Vizinhança:} Deslizar apenas uma peça verticalmente em sua coluna. Como existem 8 colunas e 7 posições de linha restantes por coluna, o fator de ramificação exato é $8 \times 7 = 56$ vizinhos iminentes.
    \item \textbf{Métrica de Qualidade ($h$):} Somatória total de conflitos. É contabilizado +1 para cada par de rainhas que divide a mesma linha ou diagonal (colunas já são garantidas como únicas pelo design do array).
    \item \textbf{Término:} Parada por Vitória ($h(s)=0$) ou Parada por Estagnação (quando todos os 56 vizinhos geram um $h \ge$ atual, caracterizando armadilhas topológicas).
    \item \textbf{Visualização do Espaço:} Imagine um terreno montanhoso formado por 16.7 milhões de arranjos distintos. O algoritmo é um alpinista vendado tentando descer aos vales ($h=0$). Depressões rasas ou superfícies planas sem rampa de descida são, respectivamente, "mínimos locais" e "platôs".
\end{itemize}

\subsection*{B) Rastreador Hill-Climbing (Steepest Ascent)}

Iniciando do tabuleiro horizontal plano $[0, 0, 0, 0, 0, 0, 0, 0]$. O escalador rigoroso varre as 56 opções e migra sempre para o estado de menor $h$.

\begin{table}[H]
\centering
\begin{tabular}{|c|c|c|}
\hline
\textbf{Salto} & \textbf{Configuração da Matriz} & \textbf{Conflitos Restantes} \\ \hline
Início & [0, 0, 0, 0, 0, 0, 0, 0] & 28 \\ \hline
1 & [0, 7, 0, 0, 0, 0, 0, 0] & 21 \\ \hline
2 & [1, 7, 0, 0, 0, 0, 0, 0] & 15 \\ \hline
3 & [1, 7, 0, 6, 0, 0, 0, 0] & 10 \\ \hline
4 & [1, 7, 0, 6, 0, 5, 0, 0] & 6 \\ \hline
5 & [1, 7, 0, 6, 0, 5, 0, 4] & 3 \\ \hline
6 & [1, 7, 0, 6, 3, 5, 0, 4] & 1 \\ \hline
7 (Preso) & [1, 7, 0, 6, 3, 5, 0, 4] & 1 \\ \hline
\end{tabular}
\end{table}

\textbf{Detalhamento de Alternativas Locais (Exemplo Iteração 6 para 7):}
\begin{itemize}
    \item No salto 6, o estado fixou-se em $h=1$. Na tentativa 7, das 56 ramificações testadas, a "melhor" transição era lateral (ex: $[1, 7, 2, 6, 3, 5, 0, 4]$ com $h=1$).
    \item O algoritmo travou formalmente em um \textbf{Platô}.
    \item A miopia do Hill-Climbing regular dita o abandono imediato da busca por inexistir rampa estritamente descendente, jogando fora todo o progresso sem atingir a coroa (0 conflitos).
\end{itemize}

\subsection*{C) Random Restart (Mitigação Bruta)}

Diante das falhas contínuas por mínimos locais, injetamos aleatoriedade reiniciando a busca a partir de gerações randômicas. Em 20 ciclos de testes independentes simulados:

\begin{table}[H]
\centering
\begin{tabular}{|c|l|c|c|}
\hline
\textbf{Ciclo} & \textbf{Semente de Origem (0-index)} & \textbf{Iterações Limite} & \textbf{Score (0 = Ideal)} \\ \hline
1 & [1, 0, 4, 3, 3, 2, 1, 1] & 4 & 1 \\ \hline
2 & [6, 0, 0, 1, 3, 3, 0, 3] & 5 & \textbf{0} \\ \hline
3 & [6, 3, 7, 4, 0, 2, 6, 5] & 2 & 2 \\ \hline
... & ... & ... & ... \\ \hline
20 & [5, 0, 3, 3, 0, 1, 0, 3] & 5 & 1 \\ \hline
\end{tabular}
\end{table}

Taxa de sucesso observada flutua em torno de 10 a 15\% globalmente.

\subsection*{D) Termodinâmica via Simulated Annealing}

Para infundir inteligência na fuga de depressões:
\begin{itemize}
    \item \textbf{Carga Térmica Original:} $T = 150.0$.
    \item \textbf{Taxa de Resfriamento:} Fator multiplicador de $0.98$ ( $T_{n+1} = T_n \times 0.98$).
    \item \textbf{Comportamento Matemático nos Abismos:}
    \begin{itemize}
        \item Custo piorando de 6 para 9 ($\Delta E = +3$): Chance de aceitação brutal ($e^{-3/147} \approx 98.0\%$).
        \item Custo piorando de 3 para 4 ($\Delta E = +1$): Aceite praticamente garantido aos $T=144.06$ ($99.3\%$).
        \item Custo piorando de 2 para 4 ($\Delta E = +2$): Tolerância maciça aos $T=141.17$ ($98.6\%$).
    \end{itemize}
    \item \textbf{Veredito Analítico:} A energia $T$ inflada no princípio age com total complacência para movimentos "estúpidos" (pioras), de modo que o platô encontrado na iteração 7 pelo Hill-Climbing teria sido saltado quase instantaneamente. No decorrer do resfriamento progressivo, as chances minguam, forçando a ancoragem num mínimo profundo real, entregando muito mais de $75\%$ de vitórias ($h=0$) na média.
\end{itemize}

\section*{Resolução da Questão 4: Verificação de Consistências em CSP}

\subsection*{A) Declaração Formal do Problema}

O paradigma CSP (Constraint Satisfaction Problem) exige o delineamento exato da alocação da equipe médica:

\begin{itemize}
    \item \textbf{Conjunto Alvo ($V$):} As vagas que demandam escala: $V = \{T_1, T_2, T_3, T_4, T_5, T_6\}$.
    \item \textbf{Opções Universais ($D$):} A bateria de profissionais autorizados para o cargo em cada vaga: $D_i = \{A, B, C, D\}$.
    \item \textbf{Regras Unitárias (1 nó):} Bloqueio pessoal. $T_3 \neq A$ (Médico A de folga no 3º).
    \item \textbf{Regras Binárias (2 nós):} 
    \begin{itemize}
        \item Vedação de turnos empilhados: $T_i \neq T_{i+1}$.
        \item Impedimento cruzado C-C: $\neg(T_2 = C \land T_5 = C)$.
    \end{itemize}
    \item \textbf{Regras Multi-nós (Globais):}
    \begin{itemize}
        \item Obrigatoriedade B: $B \in \{T_1, T_2\}$.
        \item Teto de Carga Horária D: Ocorrências $(T_i = D) \leq 2$.
    \end{itemize}
\end{itemize}

\textbf{Rede de Restrições Estruturais:}

\begin{center}
\begin{tikzpicture}[
    node distance=2.2cm,
    var/.style={circle, draw=black, thick, fill=gray!10, minimum size=1.1cm},
    edge/.style={-, thick}
]
    \node[var] (T1) {$T_1$};
    \node[var, right of=T1] (T2) {$T_2$};
    \node[var, right of=T2] (T3) {$T_3$};
    \node[var, right of=T3] (T4) {$T_4$};
    \node[var, right of=T4] (T5) {$T_5$};
    \node[var, right of=T5] (T6) {$T_6$};

    \draw[edge] (T1) -- (T2) node[midway, above] {\scriptsize $\neq$};
    \draw[edge] (T2) -- (T3) node[midway, above] {\scriptsize $\neq$};
    \draw[edge] (T3) -- (T4) node[midway, above] {\scriptsize $\neq$};
    \draw[edge] (T4) -- (T5) node[midway, above] {\scriptsize $\neq$};
    \draw[edge] (T5) -- (T6) node[midway, above] {\scriptsize $\neq$};
    
    \draw[dashdotted, blue, thick, bend left=45] (T2) to node[midway, above] {\scriptsize $\neg(C \land C)$} (T5);
\end{tikzpicture}
\end{center}

\subsection*{B) Rastreio de Backtracking Puro}

Rodando o DFS sem inteligência prévia, pela via lexicográfica cega:

\begin{table}[H]
\centering
\begin{tabular}{|c|c|l|}
\hline
\textbf{Step} & \textbf{Injeção} & \textbf{Diagnóstico da Consistência} \\ \hline
1 & $T_1 = A$ & Aceito \\ \hline
2 & $T_2 = A$ & Quebra de regra de empilhamento \\ \hline
3 & $T_2 = B$ & Aceito \\ \hline
4 & $T_3 = A$ & Quebra da folga do Médico A \\ \hline
5 & $T_3 = B$ & Quebra de regra de empilhamento \\ \hline
6 & $T_3 = C$ & Aceito \\ \hline
7 & $T_4 = A$ & Aceito \\ \hline
8 & $T_5 = A$ & Quebra de regra de empilhamento \\ \hline
9 & $T_5 = B$ & Aceito \\ \hline
10 & $T_6 = A$ & Consolidado com Êxito Total \\ \hline
\end{tabular}
\end{table}

Apenas pela sorte da ordenação inicial, 0 saltos reversos estruturais (backtracks pesados onde um nó inteiro falha) foram necessários. A montagem esbarrou em 4 becos rápidos e seguiu.

\subsection*{C) Aceleração Heurística (MRV e Degree)}

\begin{itemize}
    \item \textbf{Tática do Domínio Mínimo (MRV):} Vasculha quem tem o cinto de restrições mais apertado. $T_3$ possuía domínio nativo de apenas 3 profissionais devido à folga de A. Escolhê-lo primeiro corta raízes falhas gigantescas do espaço global.
    \item \textbf{Tática do Grau de Conexão (Degree):} Como desempate, pega-se o hub com mais cordas conectadas (ex: $T_2$, amarrado a $T_1, T_3, T_5$). Fixar $T_2$ estrangula fortemente o oxigênio (domínio) de todo o sistema simultaneamente, forçando fracassos e podas vitais cedo na árvore de busca.
\end{itemize}

\subsection*{D) Antecipação Lógica (Forward Checking)}

Ao fechar um nó, retiram-se as opções proibidas do radar dos nós futuros:

\begin{table}[H]
\centering
\resizebox{\textwidth}{!}{
\begin{tabular}{|l|l|}
\hline
\textbf{Commit} & \textbf{Sobras Lógicas nas Caixas Futuras} \\ \hline
$T_1 = A$ & $T_2: \{B\}$, $T_3: \{B, C, D\}$, Inalterados: $T_{4..6}$ \\ \hline
$T_2 = B$ & $T_3: \{C, D\}$, $T_4: \{A, B, C, D\}$, Inalterados: $T_{5..6}$ \\ \hline
$T_3 = C$ & $T_4: \{A, B, D\}$, Inalterados: $T_{5..6}$ \\ \hline
$T_4 = A$ & $T_5: \{B, C, D\}$, Inalterados: $T_6$ \\ \hline
$T_5 = B$ & $T_6: \{A, C, D\}$ \\ \hline
$T_6 = A$ & Escala Fechada. \\ \hline
\end{tabular}
}
\end{table}

A checagem reduziu o trajeto inteiro de 10 injeções (do Backtracking) para míseras \textbf{6 atuações perfeitas}. Todo movimento proibido havia sido apagado da tabela por clarividência sistêmica.

\subsection*{E) Saltos Otimizados via Backjumping (CBJ)}

O mecanismo de salto dirigido por conflito ignora o retorno cronológico tradicional, agindo de forma radical.

\begin{itemize}
    \item \textbf{Gestão da Culpa:} Toda variável rastreia intimamente quais vizinhos do passado amputaram seu domínio. Se o balde de um $T_5$ zerasse, ele não devolveria a bola para $T_4$ no cego. Ele verificaria seu \textit{Conflict Set} e acusaria pontualmente $T_2$, ordenando o retrocesso espacial direto para ele.
    \item \textbf{Impactos no Algoritmo Escalar:} Esta tática extirpa integralmente a exploração de galhos mortos (thrashing), pois varre as causas subjacentes dos problemas em vez de insistir em variáveis irrelevantes adjacentes que nada poderiam fazer para resgatar a consistência.
\end{itemize}

\section*{Resolução da Questão 5: Teoria dos Jogos (Minimax e Alpha-Beta)}

\subsection*{A) Fundamentos Teóricos Adversariais}

\begin{enumerate}[label=\roman*)]
    \item \textbf{Meta do Minimax:} Consiste em extrair o roteiro inquebrável para um protagonista, assegurando o ganho mínimo em face de um nêmese perfeito que quer arruiná-lo.
    \item \textbf{Dicotomia MAX/MIN:} Vértices MAX mapeiam o turno do nosso agente, filtrando o filho com pontuação mais gorda. Vértices MIN simulam o algoz adversário, punindo-nos ao escolher a rota com a pontuação mais miserável possível.
    \item \textbf{Carga de Utilidade:} A tradução quantitativa da vitória, derrota ou saldo do jogo contida estritamente nas folhas da árvore final.
    \item \textbf{Fluxo Retroativo:} Como a partida é lida do futuro para o presente, o sistema despenca até os nós terminais e canaliza as notas de volta para a cúpula, onde as hierarquias MAX/MIN engolem as piores ou melhores opções.
    \item \textbf{A Premissa da Perfeição:} É o postulado de que o Inimigo possui onisciência computacional e jamás sofrerá deslizes ou blefes falhos; ele jogará para nos liquidar.
    \item \textbf{Cirurgia Alpha-Beta:} Uma tática extrema de corte que poda galhos redundantes cujo acesso matemático já se prova impossível sem sequer lê-los, poupando uma massa colossal de memória e processamento.
\end{enumerate}

\subsection*{B) Avaliação do Eixo Minimax}

\textbf{Composição Analítica Final:} Após escavar todo o tronco C-D-E-H, a árvore destila a verdade matemática do jogo.

\begin{table}[H]
\centering
\begin{tabular}{|c|c|c|}
\hline
\textbf{Vértice} & \textbf{Polo} & \textbf{Avaliação Consolidada} \\ \hline
A & MAX & 6 \\ \hline
B & MIN & 5 \\ \hline
C & MAX & 5 \\ \hline
D & MAX & 9 \\ \hline
E & MIN & 0 \\ \hline
F & MAX & 2 \\ \hline
G & MAX & 0 \\ \hline
H & MIN & 6 \\ \hline
I & MAX & 7 \\ \hline
J & MAX & 6 \\ \hline
\end{tabular}
\end{table}

\begin{itemize}
    \item \textbf{Decisão de MAX na Raiz:} Enveredou para a ala direita, consolidando a utilidade global de \textbf{6}.
    \item \textbf{Trajeto Consagrado:} $A \rightarrow H \rightarrow J \rightarrow Folha(6)$.
\end{itemize}

\textbf{Modelagem da Árvore de Decisão:}
\textit{Caminho rei assinalado em amarelo forte (\textcolor{yellow}{Cor}).}

\begin{center}
\begin{forest}
  for tree={
    circle,
    draw,
    thick,
    minimum size=8mm,
    inner sep=1pt,
    s sep=5mm,
    l sep=10mm,
    font=\small
  }
  [$A(6)$, fill=yellow!50
    [$B(5)$
      [$C(5)$ [3] [5]]
      [$D(9)$ [6] [9]]
    ]
    [$E(0)$
      [$F(2)$ [1] [2]]
      [$G(0)$ [0] [-1]]
    ]
    [$H(6)$, fill=yellow!50
      [$I(7)$ [7] [4]]
      [$J(6)$, fill=yellow!50
        [5] 
        [6, fill=yellow!50]
      ]
    ]
  ]
\end{forest}
\end{center}

\subsection*{C) Execução da Poda Alpha-Beta}

\begin{table}[H]
\centering
\begin{tabular}{|c|c|c|c|c|}
\hline
\textbf{Ciclo} & \textbf{Avaliado} & \textbf{Apoio $\alpha$} & \textbf{Apoio $\beta$} & \textbf{Corte Executado?} \\ \hline
1 & C & 3 & $\infty$ & Não \\ \hline
2 & C & 5 & $\infty$ & Não \\ \hline
3 & B & $-\infty$ & 5 & Não \\ \hline
4 & D & 6 & 5 & \textbf{Sim} ($\alpha \ge \beta$) \\ \hline
5 & B & $-\infty$ & 5 & Não \\ \hline
6 & A & 5 & $\infty$ & Não \\ \hline
7 & F & 5 & $\infty$ & Não \\ \hline
8 & F & 5 & $\infty$ & Não \\ \hline
9 & E & 5 & 2 & \textbf{Sim} ($\alpha \ge \beta$) \\ \hline
10 & A & 5 & $\infty$ & Não \\ \hline
11 & I & 7 & $\infty$ & Não \\ \hline
12 & I & 7 & $\infty$ & Não \\ \hline
13 & H & 5 & 7 & Não \\ \hline
14 & J & 5 & 7 & Não \\ \hline
15 & J & 6 & 7 & Não \\ \hline
16 & H & 5 & 6 & Não \\ \hline
17 & A & 6 & $\infty$ & Não \\ \hline
\end{tabular}
\end{table}

A dinâmica fulminou ramos inteiros:
\begin{itemize}
    \item Em **D**, no instante que a folha 6 inflou o $\alpha$ local de D, o $\beta$ de B já estava sedimentado em 5. Ninguém cederia, e o nó virou lixo lógico.
    \item Em **E**, a mesma discrepância dizimou inteiramente o acesso a **G**, preservando o fluxo de dados em um nível de processamento de apenas 18 nós (contrastando com os 22 absolutos da busca completa).
\end{itemize}

\subsection*{D) Experimento: Distorção de Valores nas Folhas}

Refazendo o roteiro, injetamos instabilidade:
\begin{itemize}
    \item Folha 1 de F passou de 1 para \textbf{12}.
    \item Folha 2 de F passou de 2 para \textbf{18}.
    \item Folha 1 de I despencou de 7 para \textbf{-10}.
\end{itemize}

\textbf{Rastreio Dinâmico (Alpha-Beta Pós-Mutação):}

\begin{table}[H]
\centering
\begin{tabular}{|c|c|c|c|c|}
\hline
\textbf{Ciclo} & \textbf{Vértice} & \textbf{$\alpha$} & \textbf{$\beta$} & \textbf{Podou?} \\ \hline
... & ... & ... & ... & ... \\ \hline
6 & A & 5 & $\infty$ & Não \\ \hline
7 & F & 12 & $\infty$ & Não \\ \hline
8 & F & 18 & $\infty$ & Não \\ \hline
9 & E & 5 & 18 & Não \\ \hline
10 & G & 5 & $\infty$ & Não \\ \hline
11 & G & 5 & $\infty$ & Não \\ \hline
12 & E & 5 & 0 & \textbf{Sim} ($\alpha \ge \beta$) \\ \hline
13 & A & 5 & $\infty$ & Não \\ \hline
14 & I & -10 & $\infty$ & Não \\ \hline
15 & H & 5 & -10 & \textbf{Sim} ($\alpha \ge \beta$) \\ \hline
16 & A & 5 & $\infty$ & Não \\ \hline
\end{tabular}
\end{table}

\textbf{Conclusões Experimento Efeito Cascata:}
A presença venenosa da folha "-10" em I contaminou mortalmente o antigo caminho ouro de H. O adversário (MIN) em H viu o -10 e travou esse teto catastrófico de perdas para nós. Com $\beta = -10$, ocorreu uma poda massiva em H contra a segurança do $\alpha = 5$ da raiz A. Descartado o eixo H, fomos jogados à mercê da segunda melhor alternativa, alterando o desfecho mundial da partida inteira para o lucro mediano de \textbf{5} originado pelo bloco B. A heurística limitada não foi capaz de antecipar um desastre e a matemática fria salvou o jogador de perder com lucro negativo.

% \section*{Resolução da Questão 6: Busca Guiada por Monte Carlo (MCTS)}

\subsection*{A) Mecanismo e Pilares do Monte Carlo Tree Search}

\begin{enumerate}[label=\alph*)]
    \item \textbf{Seleção Estratégica:} A CPU escorrega pelas ramificações existentes avaliando matematicamente (UCB1) os pesos de risco/recompensa, até esbarrar no fim de um ramo inexplorado.
    \item \textbf{Bifurcação (Expansion):} Acopla-se uma nova possibilidade inédita de jogada nesse limite da árvore, destravando um braço cego para ser vasculhado.
    \item \textbf{Simulação Desenfreada (Rollout):} Daquele ponto em diante, lança-se uma chuva de lances fictícios quase cego (ou moderado por heurística), testando exaustivamente o tabuleiro até presenciar uma vitória de algum dos lados ou empate.
    \item \textbf{Eco do Sucesso (Backpropagation):} A mensagem do final do jogo "eco" flui ladeira acima pela árvore inteira, incrementando os termômetros de acerto ($W$) e uso ($N$) em cada pai no percurso.
    \item \textbf{Influência da Confiança ($N$):} O eixo divisor da fórmula. Nós hiper-testados exalam segurança empírica, derrubando severamente a sua pontuação inflacionada na equação de UCT para que a CPU pare de testá-los e vá olhar outros cantos obscuros.
    \item \textbf{Métricas de Retorno ($W$):} A base tangível do acerto. É o medidor pragmático de qualidade do braço (o foco puro de "explotação").
    \item \textbf{O Dilema Explotação vs Exploração:} Um braço de ferro em que a Exploração sacrifica lucros imediatos para iluminar cantos esquecidos da árvore (na caça por rotas geniais escondidas), contra a Explotação que suga avidamente o que já tem garantia provada estatística de dar certo. A constante $C$ atua como o juiz regulador das duas.
\end{enumerate}

\subsection*{B) Execução Rastreável (10 Pulsos)}

Com calibragem $C=1.4$ e lances randomizados:

\begin{table}[H]
\centering
\begin{tabular}{|c|c|c|c|c|}
\hline
\textbf{Pulso} & \textbf{Ramificação Adotada} & \textbf{Vértice Criado} & \textbf{Resultado} & \textbf{Ajuste dos Pesos} \\ \hline
1 & c1 & c1 & Empatou & N=1, W=0.5 \\ \hline
2 & c2 & c2 & Amarelo ganhou & N=1, W=0.0 \\ \hline
3 & c3 & c3 & Empatou & N=1, W=0.5 \\ \hline
4 & c4 & c4 & Vermelho massacrou & N=1, W=1.0 \\ \hline
5 & c4 $\rightarrow$ c1 & c1 & Vermelho massacrou & N=1, W=0.0* \\ \hline
6 & c1 $\rightarrow$ c1 & c1 & Empatou & N=1, W=0.5 \\ \hline
7 & c3 $\rightarrow$ c1 & c1 & Amarelo ganhou & N=1, W=1.0 \\ \hline
8 & c4 $\rightarrow$ c2 & c2 & Vermelho massacrou & N=1, W=0.0* \\ \hline
9 & c4 $\rightarrow$ c3 & c3 & Vermelho massacrou & N=1, W=0.0* \\ \hline
10 & c2 $\rightarrow$ c1 & c1 & Empatou & N=1, W=0.5 \\ \hline
\end{tabular}
\end{table}

\subsection*{C) Cômputo Rígido do UCT}

A calibragem exige a balança de exploração:
\[ Formula \,\, UCT = \left(\frac{w_i}{n_i}\right) + C \times \sqrt{\frac{\ln N}{n_i}} \]

Para um painel genérico de $N = 3 + 5 + 1 = 9$ trânsitos da raiz:

\begin{itemize}
    \item \textbf{Eixo 1 (c1):} $0.666 + (1.4 \times 0.855) = 1.864$
    \item \textbf{Eixo 2 (c2):} $0.800 + (1.4 \times 0.662) = 1.727$
    \item \textbf{Eixo 5 (c5):} $1.000 + (1.4 \times 1.482) = \textbf{3.075}$
\end{itemize}

O eixo **c5** absorveu a liderança. Apesar de c2 ter alta fatia real de vitórias brutas (explotação de $0.80$), ele foi visitado um excesso de vezes (punido pela raiz de log sobre $n_i$). c5 possui incerteza máxima, e a variável de bônus exploratória injetou gigantescos $2.075$ em sua veia, arrastando-o para as prioridades.

\subsection*{D) Experimento: Distorção da Personalidade da IA (C)}

Sintetizamos 500 ciclos para observar a polarização da máquina.

\begin{itemize}
    \item \textbf{Personalidade Focada ($C=0.2$):} Pulverização de nós foi $c1=21, c2=6, c3=19, c4=454$. Um vício absurdamente agressivo num só pilar que se provou levemente positivo nos primeiros instantes. Se fosse falso positivo, a máquina travaria rumo a um colapso medíocre sem margem para recuperar o prejuízo.
    \item \textbf{Personalidade Curiosa ($C=5.0$):} A pulverização foi radicalmente insana: $c1=111, c2=125, c3=130, c4=134$. A máquina não consegue confiar nos próprios resultados estelares, duvidando de seu aprendizado e revisitando todos os buracos continuamente, perdendo vasto poder de processamento em eixos podres por pânico exploratório.
\end{itemize}

\subsection*{E) O Peso do Rollout Heurístico}

A Injeção de "malícia lógica" no motor de Rollout:

\begin{table}[H]
\centering
\begin{tabular}{|p{3cm}|p{6cm}|p{6cm}|}
\hline
\textbf{Sintoma} & \textbf{Aleatoriedade Cega} & \textbf{Semi-Guloso Guiado} \\ \hline
\textbf{Aceleração de HW} & Suprema. Sem gargalos de "IFs" verificadores de tabuleiro, gera trilhões de nós/segundo. & Freada substancial de hardware por injetar if/else de heurísticas centrais em cada turno simulado. \\ \hline
\textbf{Inteligência Estatística} & Letárgica. Aprende a não cometer suicídio perdendo centenas de peças primeiro até que os blocos de W absorvam o desastre. & Veloz e Afiada. Cada amostra de simulação provê dados puros e altamente factíveis com lances razoáveis de combate humano, não-poluídos. \\ \hline
\end{tabular}
\end{table}


\appendix
\section{Códigos Python}

\subsection{Questão 1 - Buscas e Experimento}
%\lstinputlisting{q1/busca.py}
%\lstinputlisting{q1/experimento.py}

\subsection{Questão 2 - Buscas Informadas e Experimento}
%\lstinputlisting{q2/busca.py}
%\lstinputlisting{q2/experimento.py}

\subsection{Questão 3 - Busca Local (8-Rainhas)}
\lstinputlisting{q3/busca_local.py}

\subsection{Questão 4 - CSP (Escala Médica)}
\lstinputlisting{q4/csp_solver.py}

\subsection{Questão 5 - Minimax e Alpha-Beta}
\lstinputlisting{q5/jogos.py}

\end{document}
